%
% Segment 5 --- Viability & validity (slides 17--18). Steps back from the runs
% to ask what they cost and what they prove. No new figures: slide 17 distills
% the storage-economics table to the cost comparison; slide 18 is the
% internal-vs-external validity split. Numbers verbatim from the thesis
% (conclusion.tex, limitations.tex, discussion.tex).
%

{
  \renewcommand{\sectionsubtitle}{Economics, and the limits of the evidence}
  \section{Viability \& validity}
}


{
\setbeamerfont{itemize/enumerate body}{size=\small} % template pins lists to \normalsize
\footerinfootnotestrue % full-reference footnote -> move footer into footnote band
\begin{frame}{Not cost-competitive --- and that is the point}
  \framesubtitle{Storage economics}
  \small
  \begin{columns}[T, onlytextwidth]
    \begin{column}{0.44\textwidth}
      \textbf{Cost per TB-month}

      \smallskip
      \begin{tabular}{@{}lr@{}}
        EternalSeedBox            & \textbf{€241} \\
        Commercial object storage & €20--25 \\
        \hline
        Premium                   & \(\approx 10\times\) \\
      \end{tabular}

      \smallskip
      {\footnotesize Commercial \(=\) object-storage list
      price.\cite{cloud_storage_pricing}}

      \medskip
      {\footnotesize Break-even \textbf{€3.27 / node / day}
      (€98 per 30 d, 400 GB served).}
    \end{column}
    \begin{column}{0.52\textwidth}
      \begin{block}{A deliberate niche, not a general store}
        Each node pays \textbf{retail VPS} --- no economies of scale, so the rate
        cannot match commercial storage. The premium buys availability that
        \textbf{survives takedown} and open-source accountability. The market is
        users whose content no one else will host, funding the infrastructure
        directly.
      \end{block}
    \end{column}
  \end{columns}

  \textcolor{tud Orange}{\textbf{Competitive only where availability itself is the
  product.}}
\end{frame}
}


{
\setbeamerfont{itemize/enumerate body}{size=\small} % template pins lists to \normalsize
\begin{frame}{What the runs establish --- and what they cannot}
  \framesubtitle{Threats to validity}
  \small
  \begin{columns}[T, onlytextwidth]
    \begin{column}{0.46\textwidth}
      \begin{exampleblock}{Internal validity --- strong}
        \begin{itemize}
          \item Unmodified production code against faithful \textbf{Bitcoin} and
            \textbf{VPS-marketplace} replicas.
          \item Replication verified under \textbf{both} the closed and income
            conditions.
          \item The closed-system prediction was met.
        \end{itemize}
      \end{exampleblock}
    \end{column}
    \begin{column}{0.5\textwidth}
      \begin{alertblock}{External validity --- open}
        \begin{itemize}
          \item Node cap 60 --- host ceiling (\(\sim\)60--80 containers), not
            economics.
          \item Regtest drops \(\sim\)50\% of tx \(\to\) income set \(10\times\)
            rent \(\to\) economy inflated.
          \item Single VPS provider (SporeStack) --- no failover.
          \item Revenue mechanism not yet designed.
          \item Fixed genesis catalog.
        \end{itemize}
      \end{alertblock}
    \end{column}
  \end{columns}

  \textcolor{tud Orange}{\textbf{The mechanism is validated; viability is not ---
  by construction.}}
\end{frame}
}
